\scanpage{147}
\mspage{V.19}
% archive sheet number: 146

la restriction à $V_s$ du complexe $\underline\Phi^\bullet_{\unsure{\nu|v}}$ \inserted{à opérations $I$} sur $X_0$ déjà \inserted{essentiellement} introduit dans II 1 (du moins par ses faisceaux de cohomologie, qui y \deleted{\unsure{étaient}} sont notés $\underline\Phi_\nu^{q\unsure{?|2|\varepsilon}}$). Par application du foncteur oubli, ce triangle devient

(A')
\diagram{\begin{tikzcd}[ampersand replacement=\&,column sep=large,row sep=large]
\& \mathbb R\Gamma(V_{\unsure{\bar\eta|\bar y}},\Lambda_\nu)\arrow[dl] \& \\
\mathbb R\Gamma(V_s,\underline\Phi^\bullet_\nu)\arrow[rr] \& \& \mathbb R\Gamma(\bar V,\Lambda_\nu)\arrow[ul,"\bar j^*"]
\end{tikzcd}}.

\deleted{Appliquant à (A) le foncteur $\mathbb R\Gamma I$, on trouve un triangle exact}

\deleted{Lorsque} Supposons que
\[
\tag*{(4.4.1)}
(\Phi^i_\nu)_{\bar x}=0\qquad\text{pour }i<m,
\]
alors \deleted{la suite exacte associée à (A') donne}
\[
\tag*{(4.4.2)}
H^i(V_s,\underline\Phi^\bullet_\nu)=0\qquad\text{pour }i<m,
\]
et la suite exacte associée à (A') donne
\[
\tag*{(4.4.3)}
H^i(\bar V,\Lambda_\nu)\xrightarrow{\bar j^i}H^i(V_{\bar\eta},\Lambda_\nu)
\quad\begin{aligned}&\text{injectif pour }i\unsure{\leq|\geq} m-1\\&\text{bijectif pour }i<m-1.
\end{aligned}
\]
% margin note written sideways
\margin{ce qui peut aussi s'écrire sous la forme
\[
\tag*{(4.4.4)}
H^i_{\bar x}(\unsure{\bar X|\bar V},\Lambda_\nu)\longrightarrow(\underline\Phi^i_\nu)_{\bar x}\quad\text{bijectif si }\deleted{\unsure{\leq}}i<m\deleted{\unsure{4}},
\]
compte tenu des hom. canoniques
\[
\tag*{(4.4.5)}
H^i_{\bar x}(\bar X,\Lambda_\nu)\xleftarrow{\sim}\varinjlim_m H^i_{\unsure{x|\bar x}}(\unsure{V(m)|X(m)},\Lambda_\nu)\unsure{\}|*}\longrightarrow(\underline\Phi^i_\nu)_{\bar x}
\]
déduits des hom. $\underset{\text{(4.1.2)}}{H^i(\bar V,\Lambda_\nu)}-H^i(V_{\bar\eta},\Lambda_\nu)$ déduits de l'inclusion $V_{\bar\eta}\longrightarrow\bar V$, \deleted{\unsure{et des}} et des relations analogues définissant les $H^i_{\bar x}(\bar X,\Lambda_\nu)$ en fonction des $H^{\unsure{i|i-1}}(\bar V,\Lambda_\nu)$.}

En particulier, si (4.4.1) est vrai pour $m=+\infty$, par \unsure{précision|prévision|positivité} si $x$ est pt \deleted{\unsure{isolé}} de l'ens. des points de non lissité de $f:X\longrightarrow S$, on trouve que $\mathbb R(\Gamma,I)(V_s,\Phi^\bullet_\nu)=0$, i.e. que \deleted{$\unsure{\mathbb R\Gamma I}$}
\[
\deleted{\unsure{\mathbb Rb}}\mathbb R(\Gamma,I)(\bar V,\Lambda_\nu)\mathrel{\overset{\bar j^*}{\unsure{-|\longrightarrow}}}\mathbb R(\Gamma,I)(V_{\bar\eta},\Lambda_\nu)
\]
est un isom., i.e.
\[
H^i(\bar V,\Lambda_\nu)\xrightarrow[\sim]{\bar j^*}H^i(V_{\bar\eta},\Lambda_\nu),\qquad\text{pour }\forall i,
\]
\unsure{nous avons} (4.4.5) un isom pour tout $i$
